\section{The design construction}

Let $\mathcal S=S(t,k,v)$ be a finite Steiner system on a set $\Omega$:
every $t$-subset lies in exactly one block, and every block has $k$
points. We assume $2\leq t\leq k\leq v$ and write
\[
 r=\frac{v-t+1}{k-t+1}.
\]
This is the number of blocks through a fixed $(t-1)$-subset. Indeed,
their complements to that subset partition the remaining points into
sets of size $k-t+1$. Distinct blocks through that subset intersect
precisely in it. These elementary observations will account for the
exceptional incidences in the construction.

We use Construction 2 of Leemans [VN], with its attribution to Van
Nypelseer and Meixner. The incidence structure $\Gamma(\mathcal S)$ has
ordinary objects of types $1,\ldots,t-2$, namely the subsets of those
sizes. Its three remaining types are marked objects
\[
 m=(T,E),\qquad |T|=t-1,\quad T\cap E=\varnothing,
 \quad T\cup E\text{ a block},
\]
blocks $B_0$, and copied points $z\in\Omega$. Ordinary objects of
different types are incident by inclusion. Their incidences with the
last three types, and the incidences between those types, are
\[
\begin{array}{lll}
 X\sim(T,E)\iff X\subseteq T,&
 X\sim B_0\iff X\subseteq B_0,&
 X\sim z\iff z\notin X,\\
 (T,E)\sim B_0\iff B_0\cap(T\cup E)=T,&
 (T,E)\sim z\iff z\in E,&
 B_0\sim z\iff z\notin B_0.
\end{array}                                                    \tag{S1}
\]
Objects of the same type are incident only when equal. At $t=2$ the
ordinary part is empty. A marked object is then an original incident
point--line pair, so (S1) is exactly the rank-three Construction 3.

\paragraph{Theorem S.1 (chambers and firmness).}
For a nontrivial Steiner system with $t\geq2$, the construction is a
firm geometry if and only if
\[
 k\geq t+1,\qquad r\geq3.                                    \tag{S2}
\]
Under these conditions its panels missing an ordinary type have two
completions. The panels missing the marked, block and copied-point
types have respectively $k-t+2$, $r-1$ and $k-t+1$ completions.

\emph{Proof.}
The ordinary objects in a flag form a chain. Suppose first that a
marked object $(T,E)$ is present. Complete that chain inside $T$, choose
an alternative block through $T$ if none is present, and choose a copied
point of $E$ if necessary. Distinct blocks through $T$ have disjoint
tails, so these choices respect any already specified exceptional
objects and give a chamber.

If a block $B_0$ is present but no marked object is specified, extend
the largest ordinary set to a $(t-1)$-subset $T$ of $B_0$. Use the
prescribed copied point, or choose $z\notin B_0$ if none is prescribed.
The latter is possible since $r\geq3$ implies $v>k$. The unique block
$B'$ through $T\cup\{z\}$ differs from $B_0$ and meets it in $T$.
The marked object $(T,B'\setminus T)$ completes the required data.
If neither a block nor a marked object is present, use or choose a
copied point $z$ outside the largest ordinary set, extend that set to
$T\subseteq\Omega\setminus\{z\}$, and take the block on
$T\cup\{z\}$ and any alternative block through $T$. Complete the
chain inside $T$ in each case. This also treats the empty ordinary
chain when $t=2$, and proves that every flag extends to a chamber.

For a missing ordinary type the bounding subsets differ by two
elements, giving exactly two intermediate subsets. For a missing
marked type let $R$ be the largest ordinary set, of size $t-2$; put
$R=\varnothing$ if $t=2$. Choose $T=R\cup\{u\}$ with
$u\in B_0\setminus R$. The unique block through $T\cup\{z\}$
then determines the marked object, giving $k-t+2$ choices. The other
two panel counts are the number of alternative blocks through $T$
and the number of points of $E$. They are at least two under (S2).

Conversely, in a firm geometry a marked object belongs to a chamber.
Removing the copied point or the alternative block gives panels with
$k-t+1$ and $r-1$ completions. Both must have size at least two.
If $r=1$, no chamber contains that marked object in the first place.
This proves necessity as well. \hfill$\square$

In the remainder of this part we assume (S2).

\paragraph{Lemma S.2 (the induced automorphism group).}
Under (S2), the natural action identifies
$\operatorname{Aut}(\Gamma(\mathcal S))$ with
$\operatorname{Aut}(\mathcal S)$, where automorphisms preserve every
type.

\emph{Proof.}
An automorphism permutes the copied points. Their nonincidences with
blocks recover exactly each block as a subset of $\Omega$, so the
induced permutation is a design automorphism. Nonincidences with an
ordinary object recover its point subset too. For a marked object,
its copied-point neighbors recover $E$. Its block neighbors are the
$r-1\geq2$ alternative blocks through $T$; the intersection of any
two of them is $T$. Thus $T$ is recovered as well. All objects are
therefore determined by this one point permutation. Conversely a
design automorphism preserves all six rules (S1), proving the claim.
\hfill$\square$

\section{Chambers as one stabilizer fiber}

For a group acting equivariantly on a surjection $Y\to X$ between
nonempty sets, transitivity
on $Y$ is equivalent to transitivity on $X$ and transitivity of one
point stabilizer on the corresponding fiber. To prove the converse,
first move the image of a given element to the target image, and then
move within that fiber. For the forward implication, an element carrying
two points in the same fiber necessarily fixes their common image.
We include this standard action argument to make its precise use clear.

\paragraph{Theorem S.3 (the chamber fiber).}
For every $G\leq\operatorname{Aut}(\mathcal S)$ there is a
$G$-equivariant bijection
\[
\operatorname{Ch}(\Gamma(\mathcal S))\longleftrightarrow
\left\{((x_1,\ldots,x_t),B_0):
\begin{array}{l}
x_1,\ldots,x_t\text{ are distinct},\\
\{x_1,\ldots,x_{t-1}\}\subset B_0,\quad x_t\notin B_0
\end{array}\right\}.                                        \tag{S3}
\]
The fiber over each ordered $t$-tuple has size $r-1$. Consequently
$G$ is chamber-transitive if and only if it is $t$-transitive on
$\Omega$ and the stabilizer of one ordered $t$-tuple is transitive
on those $r-1$ blocks. In particular the chamber count is
\[
 (v)_t(r-1),\qquad (v)_t=v(v-1)\cdots(v-t+1).                 \tag{S4}
\]

\emph{Proof.}
The full ordinary chain and its containing $(t-1)$-set $T$ order the
points of $T$: successive differences are singletons. Read these as
$x_1,\ldots,x_{t-1}$ and take $x_t$ to be the copied point. The
alternative block is $B_0$. Conversely the data on the right of (S3)
determine the unique block $B'$ on
$T\cup\{x_t\}$. It differs from $B_0$, and the Steiner property
gives $B'\cap B_0=T$. The marked object is
$(T,B'\setminus T)$ and the ordinary chain consists of the initial
subsets of the tuple. The two constructions are inverse and commute
with the group action. Of the $r$ blocks through $T$, exactly one
contains $x_t$. The preceding stabilizer argument and elementary
counting now give all assertions. \hfill$\square$

This also proves the exact formulation of Property $(*)$ in
[VN, Theorem 9]. The source's Construction 1 has ordinary subsets of
sizes $1,\ldots,t-1$ and blocks, incident by inclusion. A chamber of
our geometry is equivalently a chamber of Construction 1, an alternative
block $B'$ through its final $(t-1)$-set $T$, and a point
$z\in B'\setminus T$. Forgetting $z$ and then $B'$ gives two
surjective equivariant maps. Applying the stabilizer argument at each
map says precisely that the action on Construction 1 is transitive,
its chamber stabilizer is transitive on the alternative blocks, and
the stabilizer of an alternative block is transitive on its remaining
points. This is Property $(*)$. The alternative encoding (S3) proves
that these conditions imply $t$-transitivity. Transitivity on the first
choice alone would not suffice. The numbers of objects are
\[
 n_i=\binom vi\quad(1\leq i\leq t-2),\qquad
 n_{\rm block}=\frac{\binom vt}{\binom kt},\qquad
 n_{\rm marked}=n_{\rm block}\binom{k}{t-1},\qquad
 n_{\rm copy}=v.                                             \tag{S5}
\]
The marked count chooses its block and then $T$. These formulas and
the panel counts determine numerical labels, but do not identify a
stabilizer group up to isomorphism.

\section{Derivation and all residues}

For $X\subseteq\Omega$ of size $j<t$, write $\mathcal S_X$ for
the derived system: its points are $\Omega\setminus X$, and its
blocks are $B\setminus X$ for blocks containing $X$. Its parameters
are $S(t-j,k-j,v-j)$. In particular $k-t+1$ and $r$ are unchanged.
Let $\mathcal B(X)$ be the geometry of the nonempty proper subsets
of a finite set $X$, with inclusion incidence and one type per size.
A join makes every object of one factor incident with every object
of another. Empty rank-zero factors are allowed.

\paragraph{Lemma S.4 (Boolean residues and the derived residue).}
The geometry $\mathcal B(X)$ and its truncation to all subset sizes
$1,\ldots,h<|X|$ are residually connected. Under (S2), the residue
at an ordinary object $X$ of size $j$ is naturally
\[
 \Gamma(\mathcal S)_X\simeq
 \mathcal B(X)*\Gamma(\mathcal S_X).                          \tag{S6}
\]

\emph{Proof.}
A flag in a subset geometry is a chain. Its residue is the join of
subset geometries on successive set differences, with a final
truncated factor above its largest member. In a factor of rank at
least two, singleton objects are connected through two-element
objects, and every other object has a singleton neighbor. If two
positive-rank factors remain, their complete cross-incidence connects
the whole residue even when each factor has rank one. Rank-zero and
rank-one residues impose no connectivity condition. This proves the
first assertion and also proves that joins of residually connected
geometries are residually connected.

Objects below $X$ give $\mathcal B(X)$. Above $X$, send an ordinary
$Y$ to $Y\setminus X$, a marked $(T,E)$ to $(T\setminus X,E)$,
a block $B_0$ to $B_0\setminus X$, and an incident copied point
$z\notin X$ to itself. These are bijections onto the objects of
$\Gamma(\mathcal S_X)$. Subtracting $X$ preserves inclusions and
all the exceptional rules (S1). Every object below $X$ is incident
with every object of this second factor, including its copied points.
Thus (S6) is an incidence isomorphism, not just an equality of counts.
\hfill$\square$

\paragraph{Theorem S.5 (reduction to rank three).}
Under (S2), $\Gamma(\mathcal S)$ is residually connected if and
only if $\Gamma(\mathcal S_R)$ is residually connected for every
$(t-2)$-subset $R$. At $t=2$ this means the identity assertion with
$R=\varnothing$. No transitivity assumption is required.

\emph{Proof.}
For necessity take a complete ordinary chain with largest member
$R$, using every size from $1$ to $t-2$. Its residue is exactly
$\Gamma(\mathcal S_R)$: the complete chain removes the Boolean
factors in (S6). Every flag residue of a residually connected
geometry is residually connected. Connectedness of the join at $R$
alone would not prove this implication, so the complete chain matters.

For sufficiency induct on $t$. The starting case $t=2$ is tautological.
For an ordinary $X$ of size $j$, the final two-design derivatives of
$\mathcal S_X$ are precisely the $\mathcal S_R$ with
$R\supseteq X$, $|R|=t-2$. Induction and Lemma S.4 imply that the
residue at $X$ is residually connected. This deals with every flag
containing an ordinary object, including flags with additional
exceptional objects.

It remains to inspect the subsets of the three exceptional types.
At a marked object $(T,E)$ the residue is the join of
$\mathcal B(T)$, the $r-1$ alternative blocks through $T$, and
the discrete set $E$ of copied points. All cross-incidences are
present. Fixing an additional block or copied point removes its
discrete factor, so Lemma S.4 treats all four cases containing the
marked object whenever the remaining rank is at least two.

At an incident block--copied-point pair $(B_0,z)$, the marked objects
are in bijection with the $(t-1)$-subsets $T$ of $B_0$: use the unique
block through $T\cup\{z\}$. The remaining incidence is precisely
the subset geometry on $B_0$ with sizes $1,\ldots,t-1$, which is
connected through points and two-subsets for $t\geq3$.

At a block $B_0$ alone, its ordinary singleton points and the copied
points outside it form a nonempty complete bipartite graph. Every
other ordinary or marked object has a singleton neighbor in $B_0$.
At a copied point $z$ alone, the ordinary and marked objects give
all subsets of $\Omega\setminus\{z\}$ of sizes $1,\ldots,t-1$;
the largest subsets again determine their unique blocks through $z$.
This connected inclusion geometry has a point neighbor on every
remaining block. Finally the full graph contains the graph between
ordinary singletons and copied points with adjacency $x\ne z$.
It is $K_{v,v}$ minus its matching, connected since $v\geq3$;
every other object has a singleton neighbor. These are the remaining
four cases. Together they exhaust the exceptional-only flags and
complete the induction. \hfill$\square$

\paragraph{Theorem S.6 (rank-two intersections).}
Under (S2), every rank-two residue of $\Gamma(\mathcal S)$ is a
partial linear space or a generalized digon. Thus the source's
property $(\mathrm{IP})_2$ holds throughout this family.

\emph{Proof.}
At $t=2$, the residue at an original flag $(x,M)$ is complete
bipartite between $M\setminus\{x\}$ and the other lines through
$x$. At an original point $p$, a remaining flag is determined by
its point $x$, since its line is $px$; incidence with the lines
avoiding $p$ is $x\in L$. This is a subincidence structure of the
original linear space. At an original line $L$, a remaining flag
$(x,M)$ is determined by $M$, since $x=M\cap L$; its incidence
with points outside $L$ is again a subincidence structure of the
original linear space. The latter two residues are partial linear
spaces.

For higher $t$, a defining flag containing an ordinary object is
handled by (S6). Residues of a join are joins of factor residues.
If two positive-rank factors remain they both have rank one and
give a generalized digon. If only the derived factor remains use
induction. In a Boolean factor, a rank-two residue is either a
join of two rank-one intervals or the incidence of singletons and
two-subsets of a three-element set. These give the same two possibilities.

A defining flag containing only exceptional objects has size $t-1$,
so here $t\leq4$. For $t=4$ all three exceptional objects are
fixed and the remaining objects are the nonempty proper subsets
of their three-set $T$, giving a triangle. For $t=3$, a marked
object and a block leave the two points of $T$ and the copied points
of $E$ with complete incidence. A marked object and a copied point
similarly leave the two points of $T$ and the alternative blocks.
A block $B_0$ and copied point $z$ leave the points and two-subsets
of $B_0$; each two-subset determines its marked object through $z$.
This is a partial linear space. These cases exhaust all possibilities.
\hfill$\square$

The Boolean argument uses the complete-chain interval decomposition.
It does not assert that arbitrary two-type truncations of a Boolean
lattice are partial linear spaces. The source constructions, firmness
criterion and transitivity target are retained from [VN]; the residue
and fiber proofs above reuse the accepted phase182/186 proof and its
independent phase187/199 reviews. The all-rank induction and its
applications are written mathematics, not a claimed full Lean theorem.

\paragraph{Reference.}
[VN] D. Leemans, \emph{A generalization of a construction due to Van
Nypelseer}, Beitr\"age zur Algebra und Geometrie 46 (2005), 561--574,
Constructions 2--3 and Theorems 1 and 9--11. The two retained
classification inputs used in the companion application are stated
there with their exact hypotheses; they do not follow from S.5.
